\section{Modelagem Fuzzy e Construção do Índice de Salvaguarda Biocultural (ISB)}

Este módulo especifica o motor de inferência fuzzy (funções de pertinência, regras e defuzzificação) para gerar o Índice de Salvaguarda Biocultural (ISB) como saída quantitativa que classifica saberes e práticas tradicionais segundo urgência de proteção e documentação.

\subsection{Fundamentação Metodológica}

A Teoria dos Conjuntos Difusos, proposta por \citeonline{Zadeh1965}, permite representar e processar informações incertas, ambíguas e subjetivas mediante atribuição de graus de pertinência a elementos em relação a conjuntos. Em contraste com a lógica booleana (pertinência 0 ou 1), a lógica fuzzy admite pertinência parcial ($\mu \in [0,1]$), possibilitando modelar conceitos linguísticos como ``alto risco de erosão intergeracional'' ou ``complexidade biocultural muito elevada''.

Para a salvaguarda de ativos tradicionais, onde as dimensões de vulnerabilidade envolvem percepções qualitativas e fronteiras imprecisas, a lógica fuzzy oferece vantagens sobre classificações lineares ou ordinais ao capturar não linearidades na interação entre fatores de risco, representar limiares de urgência e incorporar conhecimento especializado e comunitário na forma de regras linguísticas auditáveis e rastreáveis.

\subsection{Arquitetura do Sistema Fuzzy}

O sistema fuzzy proposto é do tipo Mamdani \cite{Mamdani1975}, estruturado em quatro componentes.

\subsubsection{1. Variáveis de Entrada (Antecedentes)}

Com base no painel Delphi e na adaptação WOCAT-SLM-QBR, serão definidas 5--8 variáveis de entrada linguísticas orientadas à avaliação de vulnerabilidade e prioridade de salvaguarda. As variáveis candidatas incluem Risco de Erosão Intergeracional (proporção de jovens que dominam o saber e status de transmissão ativa), Complexidade Biocultural (número de interações solo-planta-clima codificadas no saber), Singularidade Territorial (exclusividade comunitária e vínculo cosmológico local), Status de Documentação (grau de registro formal existente), Vulnerabilidade Jurídica (exposição a biopirataria e ausência de proteção formal) e Organização Social para Transmissão (estruturas comunitárias ativas de ensino e prática). Cada variável será associada a funções de pertinência triangulares ou trapezoidais, definindo termos linguísticos (Baixa, Média, Alta) com suporte em intervalos numéricos calibrados com o painel.

\textbf{Exemplo de Fuzzificação -- Risco de Erosão Intergeracional:}
\begin{itemize}
\item \textbf{Baixo:} $\mu_{Baixo}(x) = \text{trimf}(x; [0, 0, 4])$
\item \textbf{Médio:} $\mu_{Medio}(x) = \text{trimf}(x; [3, 5, 7])$
\item \textbf{Alto:} $\mu_{Alto}(x) = \text{trimf}(x; [6, 10, 10])$
\end{itemize}

\subsubsection{2. Base de Regras (Inferência)}

Serão construídas regras do tipo SE-ENTÃO que representem o conhecimento especializado consensuado pelo painel e a percepção comunitária de urgência. As regras terão a forma:

\textbf{SE} Risco de Erosão é Alto \textbf{E} Status de Documentação é Baixo \textbf{ENTÃO} ISB é Crítico

\textbf{SE} Singularidade é Alta \textbf{E} Vulnerabilidade Jurídica é Alta \textbf{ENTÃO} ISB é Muito Alto

\textbf{SE} Risco de Erosão é Baixo \textbf{E} Documentação é Alta \textbf{ENTÃO} ISB é Consolidado

A base de regras será construída de forma colaborativa em oficinas com o painel Delphi e representantes comunitários, explorando cenários extremos e intermediários. Estima-se um conjunto de 50--100 regras cobrindo as principais combinações de estados das variáveis.

\subsubsection{3. Mecanismo de Inferência}

Aplicação do operador de agregação (mínimo para ``E'', máximo para ``OU'') e cálculo do grau de ativação de cada regra. As saídas parciais serão agregadas mediante operador de união (máximo), gerando uma superfície de pertinência fuzzy para a variável de saída (ISB).

\subsubsection{4. Defuzzificação}

Conversão da saída fuzzy em valor numérico mediante método do centróide (centro de área), conforme equação:

\begin{equation}
ISB = \frac{\int \mu_{ISB}(y) \cdot y \, dy}{\int \mu_{ISB}(y) \, dy}
\end{equation}

O ISB resultante será normalizado em escala 0--100, permitindo classificação em categorias de prioridade de salvaguarda: Consolidado [0--20], Monitoramento [20--40], Atenção [40--60], Alto [60--80] e Crítico [80--100].

\subsection{Implementação Computacional}

O sistema será implementado em Python utilizando a biblioteca \texttt{scikit-fuzzy} \cite{Warner2015}, com interface gráfica para entrada de dados e visualização dos resultados. Serão geradas superfícies de decisão tridimensionais ilustrando a relação entre pares de variáveis de entrada e o ISB.

\subsection{Validação do Modelo}

\begin{enumerate}
\item \textbf{Validação lógica (face validity):} apresentação do conjunto de regras ao painel Delphi e a representantes comunitários para verificação de coerência com o conhecimento especializado e a percepção êmica de urgência.

\item \textbf{Validação empírica:} aplicação do modelo a casos reais (3--5 práticas e saberes de comunidades quilombolas do Território de Identidade Semiárido Nordeste II, Bahia) e comparação do ISB com avaliações independentes de especialistas em salvaguarda (correlação de Spearman).

\item \textbf{Análise de sensibilidade:} variação sistemática de cada variável de entrada ($\pm$ 10\%, $\pm$ 20\%) e mensuração do impacto no ISB, identificando variáveis críticas e verificando robustez do modelo.

\item \textbf{Comparação com classificação linear:} construção de modelo ordinal com as mesmas variáveis de entrada e comparação da aderência às avaliações de referência, testando se a modelagem fuzzy captura nuances que classificações lineares perdem.
\end{enumerate}

O resultado será um sistema fuzzy validado, documentado e operacional, capaz de classificar saberes segundo prioridade de salvaguarda com transparência e rastreabilidade das decisões, produzindo dossiês auditáveis para instruir processos de proteção.
